\subsection{Identifying and repairing a proof gap}
\label{sec:precision}
In our first example we identify a gap in the proof of a published generation theorem and give a short repair. The repair preserves the original algorithm and theorem statement.

In \emph{Pareto-optimal Non-uniform Language Generation},
\citet[][P13]{charikar2026pareto} construct a generator whose convergence
guarantee depends on the target language but holds for every presentation
order. Their proof of Theorem~8 uses a procedure that inserts languages
into an ordered list. When a language $L$ is inserted, the procedure
considers subcollections containing $L$, drawn from the list up to $L$,
whose intersection is finite. It stores the largest such intersection
size as a score $m^\star(L)$, together with a subcollection
$\mathcal C(L)$ attaining it. If no such subcollection exists, it stores
zero and an empty subcollection. Claim~7 asserts that the stored
subcollection remains a maximizer as further languages are inserted.

While formalizing Claim~7, GenLimitLib records a counterexample. Take two disjoint infinite
languages $A$ and $B$. When $A$ is inserted first, no subcollection
containing $A$ has finite intersection, so the procedure stores
$\mathcal C(A)=\varnothing$ and $m^\star(A)=0$. When $B$ is inserted,
$\{A,B\}$ has intersection size zero, and the insertion rule places
$B$ before $A$. However, the stored subcollection $\mathcal C(A)$
remains empty. It does not contain $A$, so it cannot be the
maximizing subcollection asserted by Claim~7.

Our repair proves that every subcollection containing $L$, drawn from the current list up to $L$ and having finite intersection, has
intersection size at most the stored score $m^\star(L)$.
It proves that this bound is preserved after each insertion.
This is the bound needed to complete the generation proof.

\subsection{Connecting papers to derive new results}
\label{sec:connections}

GenLimitLib also supports transferring constructions between different learning
problems. We connect a bounded-memory example from P31 to the replay
model of P22, obtaining a stronger impossibility example. This connection
then motivates an exact characterization result.

\emph{Language Generation with Replay} \citep[][P22]{replay}
studies a model in which the adversary can use the learner's earlier outputs to contaminate future observations. The paper studies several generation guarantees in this setting. We focus on \emph{proper generation in the limit}: proposals must come from a fixed family and eventually be contained in the true target.

A different paper, \emph{On Language Generation in the Limit with
Bounded Memory} \citep[][P31]{memory}, asks how restricting a
learner's memory affects generation. It gives a three-language
example on which an incremental index-based generator, retaining only its
previous proposed index between observations, cannot succeed. The languages are
$T\cup\{a,b\}$, $T\cup\{a,c\}$, and $T\cup\{b,c\}$,
where $T$ is infinite and $a,b,c$ are distinct points outside $T$.

Using GenLimitLib, we show that this same example also makes deterministic proper
generation under replay impossible, even with unrestricted memory.
Our proof combines P31's construction with the indistinguishability
argument in P22's impossibility proof. This improves
P22's four-language impossibility example to three languages.

This connection motivates a broader question: which finite
families permit proper generation under replay?
Our Theorem~\ref{thm:proper-replay} below gives an exact characterization.
We impose no computability restriction on the generator.
\begin{result}[When proper generation under replay is possible]
\label{thm:proper-replay}
Let $L_1,\ldots,L_N$, with $N\ge1$, be a finite family of infinite countable
languages. Deterministic proper generation under replay is possible
if and only if, for every possible first observation $x$, there is
a first proposal $L_j$ with the following property.
Group the indices $i$ with $x\in L_i$ according to equality of
$L_i\setminus L_j$. For every such group $E$, some language $L_h$
from the original family satisfies
$L_h\subseteq\bigcap_{i\in E}L_i$.
\end{result}
